\section{Implementierung}
\label{sec:implementierung}

\subsection{Mess-Engine}
Die Mess-Engine implementiert drei Messarten:

\begin{itemize}
    \item \textbf{Ping-Messung}: Nutzt \texttt{InetAddress.isReachable()}
    \item \textbf{DNS-Messung}: Nutzt \texttt{InetAddress.getByName()}
    \item \textbf{HTTP-Messung}: Nutzt \texttt{java.net.http.HttpClient}
\end{itemize}

\begin{lstlisting}[language=Java, caption=Beispiel für Ping-Messung]
public Measurement measure() throws Exception {
    long start = System.nanoTime();
    boolean reachable = InetAddress.getByName(target).isReachable(3000);
    long end = System.nanoTime();
    double latency = (end - start) / 1_000_000.0;
    return new Measurement(target, latency, reachable, "PING");
}
\end{lstlisting}

\subsection{Statistische Auswertung}
Die statistische Auswertung berechnet folgende Kennzahlen:

\begin{itemize}
    \item \textbf{Durchschnittliche Latenz}: $\mu = \frac{1}{n}\sum_{i=1}^{n}x_i$
    \item \textbf{95th Percentile}: Wert, unter dem 95\% der Messungen liegen
    \item \textbf{Paketverlust}: $\frac{\text{fehlgeschlagene Messungen}}{\text{gesamte Messungen}} \times 100$
    \item \textbf{Jitter}: Durchschnitt der Differenzen zwischen aufeinanderfolgenden Messungen
\end{itemize}

\begin{lstlisting}[language=Java, caption=Beispiel für 95th Percentile-Berechnung]
// Statistiken aus der Datenbank holen
List<Double> latencies = new ArrayList<>();
// ... Daten laden ...
Collections.sort(latencies);
int p95Index = (int) Math.ceil(latencies.size() * 0.95) - 1;
double p95Latency = latencies.get(p95Index);
\end{lstlisting}

\subsection{Web-Interface}
Das Web-Interface wurde mit Javalin implementiert:

\begin{itemize}
    \item \textbf{REST-API}: Für Datenzugriff (z. B. \texttt{/api/measurements})
    \item \textbf{Live-Chart}: Aktualisiert sich alle 5 Sekunden
    \item \textbf{Heatmap}: Zeigt Latenzentwicklung pro Stunde
    \item \textbf{Statistik-Panel}: Sofortige Einschätzung der Netzwerkqualität
\end{itemize}

\begin{figure}[H]
    \centering
    \includegraphics[width=0.8\textwidth]{../diagrams/ui.png}
    \caption{Web-Interface von SignalReport}
    \label{fig:ui}
\end{figure}